% !TEX program = xelatex
% Spatial immunosuppressive microenvironment analysis of hepatocellular carcinoma (Nature-style English version)
% Compiled with XeLaTeX using the standard article class

\documentclass[12pt, a4paper]{article}

% ============================================================
% Packages
% ============================================================
\usepackage{geometry}
\geometry{left=2.5cm, right=2.5cm, top=2.5cm, bottom=2.5cm}

\usepackage{graphicx}
\usepackage{booktabs}
\usepackage{longtable}
\usepackage{amsmath}
\usepackage{amssymb}
\usepackage{hyperref}
\usepackage{xcolor}
\usepackage{listings}
\usepackage{float}
\usepackage{caption}
\usepackage{subcaption}
\usepackage{enumitem}
\usepackage{setspace}
\usepackage{fancyhdr}
\usepackage{natbib}
\usepackage{array}
\usepackage{multirow}
\usepackage{tabularx}
\usepackage{appendix}
\usepackage{microtype}

% ============================================================
% Page layout
% ============================================================
\onehalfspacing
\pagestyle{fancy}
\fancyhf{}
\fancyhead[L]{\small Spatial immunosuppressive microenvironment of hepatocellular carcinoma}
\fancyhead[R]{\small \thepage}
\renewcommand{\headrulewidth}{0.4pt}

\lstset{
    basicstyle=\small\ttfamily,
    keywordstyle=\color{blue},
    commentstyle=\color{gray},
    stringstyle=\color{red},
    breaklines=true,
    frame=single,
    numbers=left,
    numberstyle=\tiny\color{gray},
    language=Python,
    showstringspaces=false,
    xleftmargin=1em,
    framexleftmargin=1em,
    extendedchars=false,
    keepspaces=true
}

\hypersetup{
    colorlinks=true,
    linkcolor=blue,
    citecolor=blue,
    urlcolor=blue
}

% ============================================================
% Title
% ============================================================
\title{
    \vspace{-1cm}
    \textbf{Integrating Single-Cell and Spatial Transcriptomics Reveals\\
    the Spatial Architecture of the Immunosuppressive Microenvironment\\
    in Hepatocellular Carcinoma and Its Prognostic Validation}\\[0.5em]
    \large Project Technical Report
}
\author{}
\date{\today}

% ============================================================
% Document body
% ============================================================
\begin{document}

\maketitle
\thispagestyle{empty}

\begin{abstract}

Hepatocellular carcinoma (HCC) is the third leading cause of cancer-related death worldwide, yet objective response rates to immune checkpoint inhibitors remain limited to approximately 15--20\%. This bottleneck is largely attributable to a multilayered immunosuppressive network jointly established by regulatory T cells (Tregs), tumor-associated macrophages (TAMs), and cancer-associated fibroblasts (CAFs) within the tumor microenvironment~\citep{llovet2021hepatocellular,finn2020atezolizumab}. Here we integrate single-cell RNA sequencing (scRNA-seq) with 10x Genomics Visium spatial transcriptomics, using the Cell2location Bayesian deconvolution algorithm~\citep{kleshchevnikov2022cell2location} to precisely map cell-type annotations onto every spot of the spatial sections, and applying a neighborhood-composition clustering strategy built on the ``Cellular Neighborhoods'' framework of Schürch et al.~\citep{schurch2020coordinated} to identify immunosuppressive spatial niches without relying on any pre-specified threshold. Quantitative ligand-receptor communication analysis~\citep{efremova2020cellphonedb} reveals the core signaling axes underlying the Treg--TAM--CAF triad (CXCL12--CXCR4, CCL22--CCR4, TGFB1--TGFBR1, and PD-1--PD-L1); cross-sample consistency is validated on two independent Visium sections (HCC4R and CHC20); and the extracted spatial signature is projected via single-sample gene set enrichment analysis (ssGSEA) onto the TCGA-LIHC cohort (approximately 370 patients)~\citep{tomczak2015cancer} for independent prognostic validation. Together, this project establishes a complete translational analysis framework spanning fine-scale spatial dissection, mechanistic communication analysis, and clinical cohort validation.

\noindent\textbf{Keywords:}
Hepatocellular carcinoma; spatial transcriptomics; single-cell transcriptomics; immunosuppressive microenvironment; Cell2location; Cellular Neighborhoods; ligand-receptor communication; ssGSEA; survival analysis
\end{abstract}

\newpage
\tableofcontents
\newpage

% ============================================================
%  Part I: Background and Rationale
% ============================================================

\section{Introduction}

\subsection{Disease background and clinical challenges}

Hepatocellular carcinoma (HCC) is the sixth most common malignancy worldwide, with more than 900,000 new cases and approximately 830,000 deaths annually; China alone accounts for roughly 45\% of new cases globally~\citep{llovet2021hepatocellular,chen2016liver}. Because most patients are diagnosed at an advanced stage, the five-year survival rate for HCC remains below 20\%. Sorafenib was long the standard first-line regimen for advanced HCC, yet it extended median overall survival by only about three months relative to placebo~\citep{llovet2008sorafenib}. The 2020 IMbrave150 trial demonstrated that atezolizumab combined with bevacizumab raised the objective response rate to 27.3\% and, for the first time, outperformed sorafenib in overall survival~\citep{finn2020atezolizumab}; nevertheless, approximately 70\% of patients still fail to respond to this combination, underscoring the presence of a robust, multilayered immunosuppressive mechanism within the tumor microenvironment.

Immunosuppression in the HCC microenvironment arises from the coordinated action of multiple cell types. Regulatory T cells (Tregs), classically defined by the CD$4^+$CD$25^+$FOXP$3^+$ phenotype, directly suppress the antitumor activity of effector T cells and NK cells by secreting IL-10, TGF-$\beta$, and IL-35, and by expressing checkpoint molecules such as CTLA-4, TIGIT, and LAG-3~\citep{sakaguchi2020regulatory}; single-cell transcriptomic studies have shown that Tregs are markedly enriched in HCC tissue, and the Treg-to-CD$8^+$ T-cell ratio is an independent prognostic factor~\citep{zheng2017landscape}. Tumor-associated macrophages (TAMs), predominantly polarized toward the M2 phenotype (CD163$^+$MRC1$^+$), actively recruit Tregs into the tumor via the CCL22--CCR4 and CXCL12--CXCR4 axes while secreting TGF-$\beta$ and upregulating PD-L1 to suppress effector T cells~\citep{biswas2010macrophage,noy2014tumor}; a bidirectional positive feedback loop between TAMs and Tregs sustains a stable local immunosuppressive ecology. Cancer-associated fibroblasts (CAFs), marked by \textit{FAP}, \textit{ACTA2}, and \textit{POSTN}, recruit immunosuppressive cells via CXCL12 secretion and remodel the extracellular matrix (ECM) to erect a physical barrier that impedes effector T-cell infiltration~\citep{kalluri2016biology,mariathasan2018tgfb}. The synergy of this Treg--TAM--CAF triangular axis constitutes a core mechanism of immunotherapy resistance in HCC and is the central biological object dissected in this project.

\subsection{Spatial transcriptomics and an integrative analytical paradigm}

Addressing the questions above requires an understanding of the spatial organization of immunosuppression. Conventional bulk RNA-seq homogenizes tissue prior to sequencing, thereby discarding positional information among cells; scRNA-seq resolves single-cell transcriptional features but inevitably disrupts spatial relationships during tissue dissociation~\citep{luecken2019current}. The 10x Genomics Visium spatial transcriptomics platform performs whole-transcriptome sequencing at each spatial spot (approximately 55\,$\mu$m in diameter) while preserving the native tissue architecture, allowing researchers to directly observe the spatial distribution of cell types~\citep{williams2022introduction}. Because each spot captures a mixed signal from 1--10 cells, quantitative deconvolution against an scRNA-seq reference is required to estimate its cellular composition.

Among current mainstream deconvolution tools, Cell2location employs a hierarchical negative-binomial Bayesian model to simultaneously estimate the absolute abundance of multiple cell types across the whole transcriptome while providing posterior uncertainty quantification~\citep{kleshchevnikov2022cell2location}; multiple independent benchmarks consistently show that it outperforms RCTD~\citep{cable2021robust}, CARD~\citep{ma2022integrative}, and Tangram~\citep{biancalani2021deep} in detecting rare cell types such as Tregs (typically $<$5\% in abundance)~\citep{li2022benchmarking,sang2023spotless,yan2023benchmarking}. The ``Cellular Neighborhoods'' concept developed by Schürch et al. in colorectal cancer, which applies unsupervised clustering to local cellular composition to reveal the differential contributions of distinct neighborhoods to antitumor immunity, offers a methodological paradigm for spatially resolved tumor microenvironment typing~\citep{schurch2020coordinated}. Building on this, projecting locally derived spatial gene signatures onto large clinical cohorts via ssGSEA~\citep{barbie2009systematic}, combined with multivariable Cox regression~\citep{cox1972regression} for independent prognostic validation, represents a classical strategy for translating spatial discovery into clinical relevance~\citep{hanzelmann2013gsva}.

Existing HCC studies generally lack integration along three dimensions: elevating the spatial co-localization of the Treg--TAM--CAF axis from a descriptive level to a mechanistic ligand-receptor communication level while preserving spatial information; adopting data-driven approaches to identify immunosuppressive spatial niches; and systematically validating spatial transcriptomic findings for prognostic relevance in independent, large-scale clinical cohorts. This project addresses these gaps by establishing a generalizable and complete translational analysis framework.

\subsection{Study objectives}

This project aims to answer four interrelated core scientific questions: (1) In HCC Visium sections, do immunosuppressive cellular signals (Treg, TAM, CAF) exhibit systematic spatial co-localization, and are they enriched at specific tumor-stroma interfaces? (2) Within these regions, does the Treg--TAM--CAF axis achieve coordinated immunosuppression through specific ligand-receptor pairs (e.g., CXCL12--CXCR4, CCL22--CCR4)? (3) Do the identified immunosuppressive spatial niche features exhibit cross-sample consistency across multiple independent sections? (4) Can the feature gene signature distilled from spatial data independently predict patient prognosis in the large-scale TCGA-LIHC cohort?

% ============================================================
%  Part II: Methods
% ============================================================

\section{Methods}

\subsection{Overall analytical workflow}

This study establishes a complete translational analysis framework spanning five stages—single-cell reference construction, spatial deconvolution, niche identification, mechanistic dissection, and cross-sample/clinical validation (Fig.~\ref{fig:workflow}). The overall pipeline begins with an scRNA-seq reference atlas, from which Cell2location Bayesian deconvolution maps cell-type abundance onto every spot of the Visium spatial sections; immunosuppressive niches are then identified via unsupervised clustering of spatial neighborhood composition, and a niche-specific feature gene signature is distilled through a two-layer differential expression screen. Building on this foundation, three analytical branches are pursued independently yet converge on the same set of niche feature genes: quantitative ligand-receptor communication analysis (mechanistic level), cross-sample reproducibility validation on the independent section CHC20 (spatial level), and ssGSEA projection with multivariable Cox survival analysis in the large TCGA-LIHC cohort (clinical level)—together forming a closed evidentiary chain of ``discovery--mechanism--validation.'' The specific methodological details of each stage are described in Sections 2.2--2.6 below.

\begin{figure}[h]
    \centering
    \includegraphics[width=1.0\textwidth]{plots/workflow_cn.pdf}
    \caption{Overall analytical workflow.}
    \label{fig:workflow}
\end{figure}

\subsection{Data sources and quality control}

This study integrates multimodal datasets comprising single-cell transcriptomics (scRNA-seq), spatial transcriptomics (ST), and a large-scale clinical bulk transcriptome cohort, constructing a translational medicine framework of ``peripheral discovery--spatial localization--clinical validation.''

The single-cell reference dataset was downloaded from the GEO database (accession GSE149614), encompassing core cell types including hepatocytes, T/NK cells, B cells, myeloid cells, and fibroblasts. Following current best practices in single-cell genomics~\citep{luecken2019current}, cells with total expression below 200 counts and genes expressed in fewer than 10 cells were strictly filtered out. Because the downstream Cell2location spatial deconvolution model is built on a Bayesian negative-binomial regression framework that internally models sequencing depth and technical batch effects, no conventional library-size normalization or $\log(1+p)$ transformation was applied during single-cell preprocessing; the raw count matrix was retained directly as model input.

The spatial transcriptomic sequencing data used in this study were generated on the 10x Genomics Visium high-throughput platform, employing a rigorous dual-center, cross-clinical-context design to establish a discovery-and-validation loop for spatial microenvironment features. The primary discovery section, HCC4R, was derived from a prospective clinical trial cohort (GSE238264); this sample was obtained from a resected HCC specimen following neoadjuvant immunotherapy with cabozantinib plus nivolumab, and exhibits a rich HCC-CAF spatial interaction landscape. The independent external validation section, CHC20, was derived from a separate spatial omics cohort focused on hepatocellular carcinoma associated with non-alcoholic steatohepatitis, in which the core focus is myeloid cell-mediated immunosuppression within tumor fibrotic foci.

Following best-practice quality-control standards for spatial transcriptomic analysis~\citep{williams2022introduction,palla2022squidpy}, both ST sections were uniformly filtered: only high-quality spots with total UMI counts between 5,000 and 35,000 and a mitochondrial gene fraction below 20\% were retained, while low-frequency genes expressed in fewer than 10 spots were excluded.

Finally, to systematically extend the spatial omics findings to an independent, large-scale clinical cohort for survival and prognostic validation, this study incorporated the TCGA-LIHC cohort (approximately 370 HCC patients). Normalized HTSeq-FPKM expression matrices, together with complete clinical follow-up and survival information, were downloaded from the GDC database via the TCGAbiolinks R package~\citep{colaprico2016tcgabiolinks}.

\subsection{Treg subset identification and Cell2location deconvolution}

T/NK cell subsets were extracted from the scRNA-seq data and subjected to high-resolution sub-clustering using the Leiden algorithm (resolution res=3.0)~\citep{traag2019leiden}; a bubble plot was used to visualize the expression distribution of \textit{CD3D}, \textit{CD4}, \textit{FOXP3}, and \textit{IL2RA} across subclusters. After confirming that a specific cluster displayed the classical CD4$^+$CD25$^+$FOXP3$^+$ Treg phenotype~\citep{sakaguchi2020regulatory}, this cluster was re-annotated as Treg to yield the final cell-type labels.

Cell2location was run in two stages~\citep{kleshchevnikov2022cell2location}. In the first stage, a RegressionModel was trained on the scRNA-seq data (raw counts) to derive the cell-type reference signature matrix, using a maximum of 250 training epochs, a batch size of 1024, and an early-stopping patience of 30. In the second stage, the reference signature matrix was used to independently run spatial modeling on each Visium section, with the key prior parameter N\_cells\_per\_location set to 30 (estimated from the approximately 55\,$\mu$m Visium spot diameter and typical hepatic tissue cell density~\citep{kleshchevnikov2022cell2location}) and detection\_alpha set to 20 (the officially recommended default), trained for a maximum of 1,000 epochs. The posterior mean abundance matrix (\texttt{means\_cell\_abundance\_w\_sf}) was taken as the estimate of the absolute cell-type abundance at each spot.

To facilitate result presentation and biological interpretation, in addition to the Cellular Neighborhoods clustering used to identify immunosuppressive niches, this study further assigned each spot an auxiliary, rule-based spatial-domain label (tumor\_core / tumor\_edge / stroma\_immune / other); this annotation is used solely for visualization and cross-method validation and does not participate in the primary niche-discovery pipeline. Drawing on the domain-partitioning strategies of Schürch et al.~\citep{schurch2020coordinated} and Bhate et al.~\citep{bhate2022tissue}, who used local cell-composition thresholds in multiplexed spatial omics to distinguish tumor parenchyma, invasive front, and stromal/immune regions, spots at or above the 75th percentile of hepatocyte proportion were designated as hepatocyte-high (hep\_high) and labeled tumor\_core; among the remaining non-hep\_high spots, those at or above the 50th percentile of the immune\_stroma\_score (see Appendix~\ref{app:scores}) were labeled stroma\_immune; a spot was reassigned to tumor\_edge if it additionally had at least one hep\_high neighbor within a defined physical radius, thereby characterizing the transitional interface between tumor parenchyma and the immune-stromal region; all remaining spots were labeled other.

\subsection{Identification of immunosuppressive spatial niches}

After normalizing the absolute abundances output by Cell2location into a proportion matrix, a $k$-nearest-neighbor ($k$-NN) graph was constructed on the tissue's geometric coordinates to represent spatial physical adjacency. For each spot $i$, the mean cell-type proportions across its $k$ nearest neighboring spots in physical space were computed to derive a neighborhood composition vector $\mathbf{n}_i$:
\begin{equation}
\mathbf{n}_i = \frac{1}{|\mathcal{N}_i|} \sum_{j \in \mathcal{N}_i} \mathbf{p}_j
\label{eq:neighborhood_vector}
\end{equation}
This study adopted the classical parameter choice of $k=15$ from Schürch et al.'s colorectal cancer study~\citep{schurch2020coordinated}. The neighborhood composition vector $\mathbf{n}_i$ was then treated as a high-dimensional phenotypic feature, on which a new $k$-NN graph was constructed in feature space. Finally, the Leiden community-detection algorithm (resolution $=0.5$) was run on this feature adjacency graph to adaptively identify spatial niche types. Compared with $K$-means, Leiden does not require the number of clusters to be pre-specified and theoretically guarantees cluster connectivity~\citep{traag2019leiden}.

For each Leiden cluster, the mean Treg proportion, myeloid proportion, fibroblast proportion, and immunosuppressive gene module score were computed across four functional dimensions; these were ranked in descending order within each dimension and summed to yield a combined rank $R_c$:
\begin{equation}
R_c = \text{rank}(\bar{p}_{\text{Treg},c}) + \text{rank}(\bar{p}_{\text{Myeloid},c})
+ \text{rank}(\bar{p}_{\text{Fibroblast},c}) + \text{rank}(\bar{s}_{\text{immune},c})
\label{eq:combined_rank}
\end{equation}
The cluster with the smallest $R_c$ was annotated as the immunosuppressive niche; this multidimensional composite score ensures that the identified niche simultaneously satisfies the coordinated enrichment of the Treg--TAM--CAF triangular axis rather than reflecting isolated high expression of a single cell type.

The immunosuppressive gene module score was defined as the mean CP10K+log1p expression of 14 literature-supported marker genes (\textit{FOXP3}, \textit{IL2RA}, \textit{CTLA4}, \textit{TIGIT}, \textit{LAG3}, \textit{PDCD1}, \textit{HAVCR2}, \textit{TGFB1}, \textit{IL10}, \textit{CXCL12}, \textit{CCL22}, \textit{TNFRSF18}, \textit{TNFRSF4}, \textit{IKZF2})~\citep{sakaguchi2020regulatory,sharma2023immunosuppressive}; auxiliary scores (Treg-like score, immune-stroma score, and niche score) were computed by Z-score standardizing and summing the corresponding component dimensions (see Appendix for details).

To verify that the choice of neighborhood parameters did not drive the conclusions, this study designed a continuous sensitivity analysis that avoids hard threshold cutoffs. Two neighborhood definitions were perturbed—$k$-NN neighbor count $k \in \{10, 15, 20\}$ and radius-multiplier $\text{radius\_multiplier} \in \{1.0, 1.25, 1.5\}$—applied to the composite niche score (the four-dimensional Z-score sum underlying Eq.~\ref{eq:combined_rank}). Using the score vector under the primary analysis parameters ($k=15$) as the reference, two complementary metrics—the Spearman rank correlation coefficient $\rho$ and the continuous weighted Jaccard similarity WJ (formulae defined in Appendix~\ref{app:sensitivity})—were jointly used to quantify the consistency between each perturbed score distribution and the reference. The joint interpretation criteria were as follows: $\rho > 0.95$ and $\text{WJ} > 0.90$ indicate extremely robust results; $\rho > 0.90$ and $\text{WJ} > 0.80$ indicate robust results with low sensitivity to parameter choice; $\rho < 0.85$ or $\text{WJ} < 0.70$ indicate that the parameter change substantially affects the results, warranting re-examination of the parameter choice~\citep{spearman1904proof}. This study adopted $\rho > 0.90$ and $\text{WJ} > 0.80$ as the criterion for parameter robustness, comparing the score vectors under six parameter perturbations (across the two neighborhood definitions) against the primary-analysis reference to assess whether the choice of $k=15$ materially affects downstream niche identification and feature gene extraction.

Having established the robustness of the neighborhood parameters, this study further used the continuous composite niche score (immunosuppressive\_niche\_score) to assign each spot a binarized high/low niche label, serving as the grouping variable for the subsequent two-layer feature gene screen (Section 2.6). Specifically, spots at or above the 90th percentile of this score across all spots were labeled niche\_high, and the remainder were labeled niche\_low.

\subsection{Two-layer screening of the feature gene signature}

Each Visium spot captures a mixed signal from approximately 5--50 cells (the cell-dilution effect), causing the raw expression of Treg marker genes such as \textit{FOXP3} and \textit{CTLA4} to be zero in the vast majority of spots. This severely dilutes the between-group mean difference ($\log_2$FC), such that conventional differential expression screening relying solely on the Wilcoxon test plus a $\log_2$FC threshold struggles to include these low-abundance yet biologically critical genes among the top-N candidates. To address this, in addition to $\log_2$FC, this study introduced the detection fraction (the fraction of spots with expression $>0$) as a complementary dimension, and extracted the immunosuppressive niche feature gene signature through a two-layer screening strategy.

Layer 1 (Wilcoxon + BH-FDR + AUC + detection-fraction composite ranking): Spatially adjacent spots were first aggregated into block-level means via spatial grid aggregation (block\_mwu mode), reducing the violation of test independence caused by spatial autocorrelation~\citep{palla2022squidpy}. Based on the 90th-percentile threshold of the composite immunosuppressive niche score (the four-dimensional Z-score sum underlying Eq.~\ref{eq:combined_rank}), blocks were divided into niche\_high ($\text{score} \geq P_{90}$) and niche\_low ($\text{score} < P_{90}$) groups. For each gene, a one-sided Mann-Whitney U test (i.e., the Wilcoxon rank-sum test) was performed, and $p$-values across all genes were corrected for multiple testing via the Benjamini-Hochberg FDR procedure~\citep{mann1947test}. The detection fractions of the two groups, $\text{frac}_{\text{high}}$ and $\text{frac}_{\text{low}}$, and their difference $\Delta\text{frac} = \text{frac}_{\text{high}} - \text{frac}_{\text{low}}$, were then computed, together with the AUC standardized from the Mann-Whitney U statistic~\citep{hanley1982meaning}. The primary screening criterion was FDR $<$ 0.1, AUC $>$ 0.60, and ($\log_2$FC $>$ 0.3 OR $\Delta$frac $>$ 0.10); if fewer than 80 genes satisfied this, the criterion was relaxed to FDR $<$ 0.2, AUC $>$ 0.55, and ($\log_2$FC $>$ 0.2 OR $\Delta$frac $>$ 0.05). Candidate genes passing this screen were further filtered against a blacklist of more than 350 non-specific genes—including housekeeping genes, ribosomal proteins, mitochondrial genes, hemoglobin genes, cell-cycle genes, broad-spectrum inflammatory genes, and tissue damage/stress-response genes—to ensure niche specificity rather than general tissue activity. Candidates were finally ranked by a weighted composite score,
\begin{equation*}
\text{signature\_rank\_score} = 0.35 \cdot \text{minmax}(\log_2\text{FC}) + 0.35 \cdot \text{minmax}(\Delta\text{frac}) + 0.30 \cdot \text{minmax}(\text{AUC})
\label{eq:signature_rank_score}
\end{equation*}
and the top 80 genes (without truncation if forced inclusion of prior genes exceeded this number) were taken as the preliminary signature gene set.

Layer 2 (targeted AUC testing and forced inclusion of prior functional gene sets): For three literature-supported prior functional gene sets—Treg markers (\textit{FOXP3}, \textit{IL2RA}, \textit{CTLA4}, etc.), TAM signature genes (\textit{CD163}, \textit{MRC1}, \textit{TGFB1}, etc.), and CAF activation genes (\textit{FAP}, \textit{ACTA2}, \textit{COL1A1}, etc.)—AUC, $\log_2$FC, the detection-fraction difference $\Delta$frac, and the FDR-corrected $p$-value were computed individually between the niche\_high and niche\_low groups to independently assess whether each prior gene exhibited niche-specific enrichment. Prior genes satisfying AUC $>$ 0.60 and $\Delta$frac $>$ 0.05 (and not present on the non-specific gene blacklist) were forcibly incorporated into the final signature gene list regardless of whether they had already passed the genome-wide ranking in Layer 1, preventing biologically well-established key genes from being overlooked owing to sparse expression in Visium data. The complete implementation of the two-layer screening pipeline is provided in Appendix~\ref{app:code}.

\subsection{Ligand-receptor communication analysis}

To elevate the descriptive cell co-localization findings above to a mechanistic communication level, this study quantitatively compared the communication strength of 12 literature-supported ligand-receptor (L-R) pairs—including \textit{CXCL12}--\textit{CXCR4}, \textit{CCL22}--\textit{CCR4}, \textit{TGFB1}--\textit{TGFBR1}, and PD-1--PD-L1 (complete list in Appendix~\ref{app:lr_table})—between the niche\_high and niche\_low groups. Because the conventional ligand-receptor product method~\citep{efremova2020cellphonedb} is prone to yielding a product of exactly zero at the Visium spot level whenever either gene's expression approaches zero (the cell-dilution effect), thereby completely masking signals from Treg-specific pathways such as \textit{CCL22}--\textit{CCR4}, this study first rank-normalized each gene's expression vector across all spots (mapping values to $[0,1]$) and then used the product of ligand and receptor ranks as a surrogate for communication strength:
\begin{equation*}
\text{score}_{LR,i} = \text{rank}(E_{L,i}) \times \text{rank}(E_{R,i}), \quad \text{rank}(\cdot) \in [0,1]
\end{equation*}
This rank-normalized product preserves the relative ordering information of ligand and receptor expression while effectively avoiding the dominance of sparse zero values over the product, and naturally satisfies the logical AND constraint that both ligand and receptor must exhibit relatively high expression. For each L-R pair, the mean rank-normalized product was computed separately within the niche\_high and niche\_low groups, and $\log_2(\text{mean}_{\text{high}}/\text{mean}_{\text{low}})$ (with a pseudocount $\varepsilon = 10^{-6}$ added to the denominator to avoid division by zero) was used to quantify the direction of enrichment or depletion within the immunosuppressive niche; the differential pattern of the 12 communication axes between the two groups was ultimately presented jointly via a heatmap (communication strength) and a bar chart ($\log_2$FC).

Beyond cell-cell communication, this study further used predefined functional gene set scores to systematically characterize the activity of immunosuppression, stromal remodeling, inflammation, and metabolic stress across spatial domains at the pathway level (Fig.~\ref{fig:domain_func}). Specifically, drawing on the construction logic of the Hallmark gene sets~\citep{liberzon2015molecular} and the general paradigm of single-sample gene set scoring~\citep{barbie2009systematic}, ten functional pathway gene sets closely related to the Treg--TAM--CAF immunosuppressive axis and its accompanying microenvironmental changes were manually curated—TGF-$\beta$ Signaling, EMT, Angiogenesis, IL6-JAK-STAT3, Treg Immune Suppression, TAM Polarization (M2), CAF Activation, Immune Checkpoint, Hypoxia, and TNF-$\alpha$ Signaling via NF-$\kappa$B—each comprising six representative genes. For each spot, the mean CP10K+log1p expression of all available genes within a pathway was computed as that pathway's single-sample score, which was then averaged within each of the four rule-based spatial-domain labels described in Section 2.5 (tumor\_core / tumor\_edge / stroma\_immune / other) to yield the mean score of each pathway across the four spatial domains. In the resulting bubble plot, bubble size represents the relative magnitude of the score, and color indicates whether the score exceeds the pathway's global mean across all four domains (red denotes above the mean, blue denotes below).

\subsection{Cross-section consistency validation and TCGA survival analysis}

Using the signature gene set extracted from HCC4R as a reference, module scores were computed on the CHC20 section separately for the immune module genes (Treg/TAM-related, 21 genes in total) and the stromal/ECM module genes (CAF/ECM-related, 21 genes in total). These were Z-score standardized and summed to obtain a combined niche score (immune\_stromal\_niche\_score), and the cross-sample consistency of the score distributions was quantified via spatial projection maps and the Mann-Whitney U test~\citep{mann1947test}.

Survival analysis employed the ssGSEA method~\citep{barbie2009systematic} implemented in the GSVA R package~\citep{hanzelmann2013gsva}, taking the $\log_2(\text{FPKM}+1)$ matrix as input with a Gaussian kernel (kcdf=Gaussian) and absolute ranking (abs.ranking=TRUE) to compute a scalar enrichment score for each patient in the TCGA-LIHC cohort. ssGSEA independently generates a continuous score reflecting niche activity for each patient, which is better suited to inclusion in a prognostic regression model than the cross-sample relative activity estimates produced by GSVA~\citep{barbie2009systematic}. A multivariable Cox proportional hazards regression model~\citep{cox1972regression} was constructed with the ssGSEA score, age, and tumor stage as covariates:
\begin{equation*}
h(t \mid \mathbf{x}) = h_0(t) \cdot \exp(\beta_1 s_{\text{ssGSEA}} + \beta_2 \text{age} + \beta_3 \text{stage})
\label{eq:cox}
\end{equation*}
After adjusting for age and stage (well-established strong independent prognostic factors in HCC~\citep{llovet2021hepatocellular}), a significant $\beta_1$ for the ssGSEA score would demonstrate that the immunosuppressive niche signature captures prognostic information independent of clinical stage. Kaplan-Meier survival curves~\citep{kaplan1958nonparametric} were plotted using the median ssGSEA score as the threshold, and the log-rank test~\citep{mantel1966evaluation} was used to assess between-group survival differences.

% ============================================================
%  Part III: Results
% ============================================================

\section{Results}

\subsection{Construction of the scRNA-seq reference atlas and identification of the Treg subset}

\begin{figure}[h]
    \centering
    \begin{subfigure}[b]{0.48\textwidth}
        \centering
        \includegraphics[width=\textwidth]{plots/cell_type_marker_gene_expression.png}
        \caption{Heatmap of marker gene expression for major cell types.}
        \label{fig:cell_type}
    \end{subfigure}
    \hfill
    \begin{subfigure}[b]{0.48\textwidth}
        \centering
        \includegraphics[width=\textwidth]{plots/treg_marker_gene_expression_t_nk_subclusters.png}
        \caption{Bubble plot of Treg marker gene expression across T/NK subclusters.}
        \label{fig:treg}
    \end{subfigure}
    \caption{Cell-type identification in the single-cell reference atlas and fine-scale identification of the regulatory T-cell (Treg) subset.}
    \label{fig:scRNA-seq_Treg}
\end{figure}

The single-cell reference atlas was constructed from the GSE149614 scRNA-seq dataset. The cell-type marker gene heatmap (Fig.~\ref{fig:cell_type}) displays the Z-score-normalized mean expression profile of each cell type, with the horizontal axis representing the seven major identified cell types (including T/NK, Treg, and myeloid cells) and the vertical axis representing the classical marker genes of each type; color encodes the relative Z-score-normalized expression level. Hepatocytes/malignant cells (\textit{ALB}, \textit{APOA2}, \textit{APOA1}), T/NK cells (\textit{CD3D}, \textit{GNLY}, \textit{GZMB}), myeloid cells (\textit{LYZ}, \textit{AIF1}, \textit{HLA-DRA}), HSCs/fibroblasts (\textit{COL1A1}, \textit{ACTA2}, \textit{RGS5}), B cells (\textit{MS4A1}, \textit{CD79A}), and endothelial cells (\textit{PECAM1}, \textit{CDH5}) all display a clear diagonal enrichment pattern, validating the reliability of the cell-type annotation~\citep{ma2019comprehensive}.

For the T/NK subset, high-resolution Leiden clustering at res=3.0 followed by a horizontal bubble plot (Fig.~\ref{fig:treg}) shows the expression of \textit{CD3D}, \textit{CD4}, \textit{FOXP3}, and \textit{IL2RA} across the different T/NK subclusters, where bubble size represents the fraction of cells expressing the gene and color intensity represents mean expression level. A specific cluster (cluster 9) exhibits the largest and darkest \textit{FOXP3} bubble, co-expression of \textit{CD4}, and \textit{IL2RA} markedly higher than in other clusters—a combined pattern consistent with the classical CD4$^+$CD25$^+$FOXP3$^+$ Treg phenotype~\citep{sakaguchi2020regulatory}; \textit{FOXP3} is nearly absent in the remaining T/NK subclusters, confirming the subcluster specificity of \textit{FOXP3} expression. Cluster 9 was re-annotated as Treg to generate the final cell-type label, which was incorporated into the reference signature for Cell2location deconvolution.

\subsection{Cell2location deconvolution reveals the spatial distribution of immune cells in HCC4R}
\subsubsection{Spatial distribution of immune cells in the HCC4R section}
\begin{figure}[h]
    \centering
    \includegraphics[width=0.6\textwidth]{plots/hcc4r_he_image.png}
    \caption{Anatomical delineation of spatial niches and the tumor microenvironment in hepatocellular carcinoma (HCC4R).}
    \label{fig:he}
\end{figure}

Figure~\ref{fig:he} precisely maps the spatial niche landscape of the HCC4R section onto every spot. The accompanying H\&E-stained image clearly reveals the morphological boundary between the tumor parenchyma (large, densely stained nuclei, hepatocyte-rich) and the stromal/interstitial region (fibrosis, inflammatory infiltration), providing a histological reference for the spatial transcriptomic results.

The spatial-domain map shows that Tumor Core (red) forms large, continuous regions that constitute the main structural backbone of the tissue section, in close spatial agreement with the nucleus-dense regions revealed by H\&E staining and the hepatocyte deconvolution abundance map; Stroma/Immune (blue) is embedded as band-like structures surrounding Tumor Core and along the tissue margin, forming a continuous region that encircles and interdigitates with the tumor parenchyma rather than appearing as isolated scattered points—consistent with the methodological expectation that unsupervised clustering effectively captures continuous tissue spatial microstructure rather than random noise~\citep{schurch2020coordinated}; Tumor Edge (orange) is distributed along the transitional zone between Tumor Core and Stroma/Immune, forming a narrow band-like interface between the two and occupying the smallest proportion of the section. This spatial pattern suggests that the tissue architecture of the HCC4R section is not a strict dichotomy between tumor and stroma, but rather exhibits a continuous ``core--edge--stromal/immune'' transitional gradient, consistent with the classical spatial description of the ``immune-excluded'' HCC phenotype, in which tumor cells occupy the core while immune-stromal cells are confined to the peripheral boundary~\citep{llovet2021hepatocellular,binnewies2018understanding}.

\begin{figure}[h]
    \centering
    \begin{subfigure}[b]{0.48\textwidth}
        \centering
        \includegraphics[width=\textwidth]{plots/spatial_hepatocyte.png}
        \caption{Spatial proportion distribution of hepatocytes.}
        \label{fig:spatial_hepatocyte}
    \end{subfigure}
    \hfill
    \begin{subfigure}[b]{0.48\textwidth}
        \centering
        \includegraphics[width=\textwidth]{plots/spatial_treg.png}
        \caption{Spatial abundance distribution of Treg-like cells.}
        \label{fig:spatial_treg}
    \end{subfigure}

    \vspace{0.3cm}

    \begin{subfigure}[b]{0.48\textwidth}
        \centering
        \includegraphics[width=\textwidth]{plots/spatial_myeloid.png}
        \caption{Spatial proportion distribution of myeloid cells.}
        \label{fig:spatial_myeloid}
    \end{subfigure}
    \hfill
    \begin{subfigure}[b]{0.48\textwidth}
        \centering
        \includegraphics[width=\textwidth]{plots/spatial_fibroblast.png}
        \caption{Spatial proportion distribution of fibroblasts.}
        \label{fig:spatial_fibroblast}
    \end{subfigure}
    \caption{Spatial abundance and proportion distribution of major cell types mapped by Bayesian deconvolution.}
    \label{fig:spatial_cell_abundance}
\end{figure}

The spatial cell-type distribution maps (Fig.~\ref{fig:spatial_cell_abundance}) reveal a spatial pattern highly consistent with expectations. Hepatocyte-high regions are concentrated within the tumor parenchyma, spatially matching the nucleus-dense regions in the H\&E image (Fig.~\ref{fig:spatial_hepatocyte}); Tregs display a focal, right-skewed enrichment pattern, with high-abundance hotspots predominantly localized at the interface between tumor parenchyma and stroma (Fig.~\ref{fig:spatial_treg}), consistent with the biological expectation that Tregs preferentially accumulate at the ``tumor-immune interface''~\citep{zheng2017landscape}; myeloid cells show a marked spatial co-localization trend with Tregs (Fig.~\ref{fig:spatial_myeloid}), echoing the signaling mechanism by which TAMs actively recruit Tregs via the CCL22--CCR4 axis~\citep{noy2014tumor}; fibroblasts are predominantly distributed within the stromal layer surrounding the tumor parenchyma (Fig.~\ref{fig:spatial_fibroblast}), reflecting the functional localization of CAFs in constructing a physical barrier~\citep{kalluri2016biology}.

\begin{figure}[h]
    \centering
    \includegraphics[width=0.8\textwidth]{plots/cell_type_spatial_colocalization.png}
    \caption{Heatmap of spatial co-localization correlations among major cell types.}
    \label{fig:hot_plot}
\end{figure}

Pearson correlation analysis among cell types (Fig.~\ref{fig:hot_plot}) further quantitatively characterizes the spatial co-localization relationships described above. The correlation heatmap shows that myeloid cells exhibit moderate, significant positive correlations with both Tregs ($r = 0.35$) and fibroblasts ($r = 0.35$), whereas fibroblasts and Tregs show a weak negative correlation ($r = -0.15$). This ``asymmetric triangular'' correlation pattern suggests that the immune-stromal compartments may be subject to fine-grained functional stratification in space: myeloid cells (TAMs) act as a central spatial hub that associates simultaneously with both immunosuppressive and stromal cells—on the one hand recruiting and anchoring Tregs by secreting chemokines such as CCL22 and CXCL12~\citep{noy2014tumor,curiel2004specific}, and on the other hand cooperating with fibroblasts (CAFs) in extracellular matrix remodeling and physical barrier construction~\citep{kalluri2016biology}—whereas the direct spatial association between fibroblasts and Tregs is weak, consistent with the hypothesis that these two cell types are more likely coupled indirectly through the shared TAM node rather than through direct mutual recruitment; this aligns with prior observations that TAMs act as a ``signaling hub'' within the tumor immunosuppressive network~\citep{biswas2010macrophage}.

Furthermore, all three immune-stromal compartments (myeloid, Treg, fibroblast) exhibit strong negative spatial correlations with the parenchymal hepatocyte compartment (correlation coefficients of $r = -0.73$, $r = -0.62$, and $r = -0.55$, respectively), indicating that these three compartments are spatially concentrated within non-parenchymal regions and systematically excluded from the tumor core. This quantitative evidence strongly supports a spatial architecture in which the immunosuppressive niche is chiefly entrenched at the tumor-stroma transitional zone rather than diffusely distributed throughout the section—a topological feature that impedes effective infiltration of effector immune cells into the tumor core, closely matching the classical spatial description of ``immune-excluded'' tumors, in which immune cells are physically confined to the tumor boundary and stromal region and are unable to penetrate the tumor parenchyma to exert cytotoxic function~\citep{llovet2021hepatocellular,chen2017oncology}.

\subsubsection{Differences among spatial domains in HCC4R}

\begin{figure}[h]
    \centering
    \begin{subfigure}[b]{0.8\textwidth}
        \centering
        \includegraphics[width=\textwidth]{plots/cell_type_composition_per_spatial_domain.png}
        \caption{Cell-type composition of the spatial domains.}
        \label{fig:cell_type_composition}
    \end{subfigure}
    \vspace{0.5cm}

    \begin{subfigure}[b]{0.32\textwidth}
        \centering
        \includegraphics[width=\textwidth]{plots/treg_like_score_by_spatial_domain.png}
        \caption{Treg-like score.}
        \label{fig:treg_score}
    \end{subfigure}
    \hfill
    \begin{subfigure}[b]{0.32\textwidth}
        \centering
        \includegraphics[width=\textwidth]{plots/immune_stroma_score_by_spatial_domain.png}
        \caption{Immune-stroma score.}
        \label{fig:immune_stroma}
    \end{subfigure}
    \hfill
    \begin{subfigure}[b]{0.32\textwidth}
        \centering
        \includegraphics[width=\textwidth]{plots/immunosuppressive_niche_score_by_spatial_domain.png}
        \caption{Immunosuppressive niche score.}
        \label{fig:immunosuppressive_niche}
    \end{subfigure}
    \caption{Cell-type composition and module scores across different spatial domains.}
    \label{fig:spatial_domain_scores}
\end{figure}

The stacked bar chart of cellular composition (Fig.~\ref{fig:cell_type_composition}) provides quantitative support for the spatial pattern described above. The hepatocyte proportion shows a clear decreasing gradient across the four domains: Tumor Core (approximately 35\%) $>$ Other (approximately 24\%) $>$ Tumor Edge (approximately 25\%) $>$ Stroma/Immune (lowest), consistent with the definition of Tumor Core as dominated by hepatocytes/malignant cells. Conversely, the proportions of Treg, myeloid, and fibroblast cells—the three core cell types of the immunosuppressive axis—are all higher in Stroma/Immune and Tumor Edge than in Tumor Core: the Treg proportion reaches its highest level among the four groups in Stroma/Immune, and myeloid and fibroblast cells show similarly enriched trends; the coordinated enrichment of all three within the same spatial domain directly confirms the spatial synergy of the Treg--TAM--CAF triangular immunosuppressive axis described in the main text~\citep{sakaguchi2020regulatory,noy2014tumor,kalluri2016biology}. The T/NK cell proportion is relatively prominent in Tumor Edge, consistent with the notion that the tumor invasive front is the primary battleground where effector T cells confront regulatory cells~\citep{schurch2020coordinated}. The endothelial cell proportion remains relatively stable across the four groups, suggesting that vascular distribution per se is not a dominant driver of this clustering-based partition.

The violin plots of module scores across spatial domains (Figs.~\ref{fig:treg_score}, \ref{fig:immune_stroma}, and \ref{fig:immunosuppressive_niche}) yield similar conclusions. Quantitative comparison across spatial domains reveals significant differences in microenvironmental features among regions. Kruskal-Wallis tests show that the overall distributions of all three scores differ highly significantly among regions (all $p < 0.001$). Specifically, the immune-stroma score ($\eta^2 = 0.83$), the immunosuppressive niche score ($\eta^2 = 0.41$), and the Treg-like score ($\eta^2 = 0.33$) in the Stroma/Immune region are all significantly higher than in Tumor Core, Tumor Edge, and Other (all pairwise comparisons $p < 0.001$). This difference is most pronounced for the immune-stroma score: non-Stroma/Immune regions are heavily concentrated in the low-score range, whereas the Stroma/Immune region exhibits a pronounced right-tail shift, spreading broadly into the high-score range~\citep{mann1947test}.

Taken together, the quantitative results from both cellular composition and functional scoring indicate that the HCC4R section exhibits a monotonic immunosuppressive gradient of ``low in Tumor Core, intermediate in Tumor Edge, and high in Stroma/Immune,'' and that this rule-based domain partitioning is statistically supported in both effect size and significance, providing a prior expectation for the data-driven Cellular Neighborhoods clustering presented in the following section.

\subsection{Cellular Neighborhoods clustering identifies the immunosuppressive spatial niche}

\begin{figure}[h]
    \centering
    \includegraphics[width=0.8\textwidth]{plots/neighborhood_composition_clusters.png}
    \caption{Leiden clustering results for the immunosuppressive spatial niche.}
    \label{fig:niche_cluster}
\end{figure}

Figure~\ref{fig:niche_cluster} shows the clustering result obtained by reconstructing a graph in feature space from the neighborhood composition vector $\mathbf{n}_i$ described in Section 2.4 ($k$-NN, $k=15$) and running Leiden clustering (resolution $= 0.5$). This figure directly colors the spatial coordinate system by cluster identity (NC-0 through NC-11, 12 neighborhood clusters in total), fully preserving the tissue heterogeneity structure natively discovered by the Leiden algorithm within neighborhood-composition feature space.

The most prominent feature in the figure is that the Immunosuppressive Niche forms the single largest and most highly connected cluster in space, occupying the central region of the section and extending in multiple directions, with a clear advantage in spatial continuity and coverage relative to the remaining 11 clusters. This cluster was determined by the composite ranking formula in the main methodology (Eq.~\ref{eq:combined_rank}) to be the immunosuppressive\_niche—that is, it has the smallest sum of descending ranks across the four dimensions of Treg proportion, myeloid proportion, fibroblast proportion, and immunosuppressive gene score, and is therefore judged to be the neighborhood type with the strongest overall immunosuppressive character under the joint weighting of four-dimensional evidence. The Immunosuppressive Niche is spatially contained within the Stroma/Immune and Tumor Edge regions of Fig.~\ref{fig:he}; the mutual corroboration of these two independent annotation schemes (native Leiden cluster identity vs. rule-based semantic labels) in spatial location methodologically strengthens the robustness of this immunosuppressive niche determination and is consistent with the expectation that the cluster map and the semantic annotation map should exhibit a high degree of spatial concordance~\citep{schurch2020coordinated,goltsev2018deep}.

Another feature clearly visible in the figure is that, aside from NC-2, the remaining clusters carry clear histological significance: for example, NC-1 and NC-3 each form independent, large, continuous regions in the lower-left and upper-right corners of the section, respectively, likely corresponding to distinct batches or regional subtypes of tumor parenchyma; NC-5 and NC-9, among others, are embedded as narrower bands between larger clusters such as NC-1 and NC-2, suggesting that these neighborhoods represent transitional substructures between tumor and stroma. Overall, the Leiden algorithm, without any prior specification of cluster number, automatically identified 12 internally consistent and spatially continuous neighborhood types, validating that unsupervised clustering can indeed capture genuine tissue spatial microstructure rather than artifacts arising from arbitrary parameter choices~\citep{traag2019leiden,schurch2020coordinated}. This pattern of ``one dominant immunosuppressive neighborhood coexisting with multiple tumor/stromal substructure neighborhoods'' is also consistent with recent observations in spatial proteomics studies that tumor tissues commonly harbor multiple coexisting cellular neighborhoods, with immunosuppressive neighborhoods typically manifesting as continuous structures spanning the tumor-stroma interface rather than isolated hotspots~\citep{jackson2020single,bhate2022tissue}, further supporting that the immunosuppressive niche identified here represents a genuine spatial entity with a stable histological basis in the HCC4R section, rather than a chance artifact of arbitrary clustering parameter choices.

\subsection{Extraction and functional annotation of the niche feature gene signature}

\begin{figure}[h]
    \centering
    \begin{subfigure}[b]{0.71\textwidth}
        \centering
        \includegraphics[width=\textwidth]{plots/fraction_of_spots_enrichment.png}
        \caption{Scatter plot of feature gene screening for the high- versus low-niche comparison.}
        \label{fig:high_low_niche_gene}
    \end{subfigure}
    \hfill
    \begin{subfigure}[b]{0.27\textwidth}
        \centering
        \includegraphics[width=\textwidth]{plots/niche_signature_genes_by_spatial_domain.png}
        \caption{Heatmap of core differentially expressed gene abundance across spatial domains.}
        \label{fig:domain_hotplot}
    \end{subfigure}
    \caption{Identification and landscape of differentially expressed genes across distinct tumor microenvironment substructures.}
    \label{fig:niche_core_gene}
\end{figure}

Figure~\ref{fig:high_low_niche_gene} presents the Layer 1 detection-fraction compensation results, with $\log_2$FC on the horizontal axis and $\Delta$fraction (the detection-fraction difference) on the vertical axis: red points (Both-sig, $n = 33$) satisfy both significance criteria simultaneously, including high-confidence immunosuppressive genes such as \textit{CXCL12} and \textit{TGFB1}, consistent with the known chemokine axis and CAF/Treg secretory mechanisms~\citep{domanska2013targeting,mariathasan2018tgfb}. Notably, the $\log_2$FC of most immunosuppressive candidate genes is concentrated between 0.1 and 0.3 despite an already high $\Delta$fraction—an apparent anomaly explained by the cell-dilution effect inherent to Visium data~\citep{luecken2019current}, underscoring the necessary compensatory role of the $\Delta$frac dimension relative to conventional $\log_2$FC screening.

Figure~\ref{fig:domain_hotplot} shows a Z-score heatmap of the expression of the top genes, ranked in descending order of composite\_score, across the four spatial domains; nearly all signature genes exhibit the consistent trend of highest expression in Stroma/Immune, lowest in Tumor Core, and intermediate in Tumor Edge, closely matching the spatial pattern of Fig.~\ref{fig:he}~\citep{zheng2017landscape}. The heatmap also shows that MHC class II molecules and lymphocyte-related genes (\textit{CD52}, \textit{TRAC}, etc.) co-enrich with stromal/fibroblast genes (\textit{DCN}, \textit{LUM}, etc.) in the Stroma/Immune domain, suggesting that this region harbors both active antigen presentation and stromal remodeling features~\citep{neefjes2011towards,iozzo2015proteoglycan}.

\begin{figure}[h]
    \centering
    \includegraphics[width=0.9\textwidth]{plots/pathway_enrichment_by_spatial_domain.png}
    \caption{Enrichment landscape of predefined functional pathway gene set scores across spatial domains.}
    \label{fig:domain_func}
\end{figure}

Figure~\ref{fig:domain_func} presents a bubble plot of the enrichment of ten predefined functional pathway gene set scores across the four spatial domains. Globally, Tumor Core exhibits the lowest activity across all ten pathways, while nine of the ten pathways (all except Hypoxia) peak in the Stroma/Immune domain, confirming this region as the core active site for pro-tumorigenic, stromal-remodeling, and immunosuppressive signaling. Regarding differential patterns, CAF Activation is enriched exclusively in Stroma/Immune, consistent with the mechanism by which CAFs secrete CXCL12 and remodel the ECM~\citep{kalluri2016biology}; EMT, TGF-$\beta$, and Treg Immune Suppression pathways are elevated in both Tumor Edge and Stroma/Immune, consistent with the expectation of coordinated activity between the invasive front and the immunosuppressive core~\citep{nieto2016emt,sakaguchi2020regulatory}. Notably, the Hypoxia pathway is instead enriched in Tumor Edge rather than Stroma/Immune, possibly reflecting active neoangiogenesis and an acute oxygen-supply mismatch in this region that more readily triggers detectable acute HIF pathway activation~\citep{palazon2017hif}.

Taken together, the gene signature and functional pathway scores extracted in this section jointly confirm the spatial localization of Stroma/Immune as the core immunosuppressive site and Tumor Edge as the invasive-front transitional zone, providing a rigorously denoised feature gene set for the subsequent cross-sample validation and survival analysis~\citep{zheng2017landscape,schurch2020coordinated}.

\subsection{Spatial reproducibility of the niche gene program in the independent section CHC20}

This section presents the cross-section consistency validation described in Section 2.7, testing whether the immunosuppressive niche signature gene set extracted from the HCC4R section exhibits cross-sample generalizability beyond a single section.

\begin{figure}[h]
    \centering
    \begin{subfigure}[b]{0.48\textwidth}
        \centering
        \includegraphics[width=\textwidth]{plots/chc20_immune_module_spatial_score.png}
        \caption{Immune module score.}
        \label{fig:immune_score}
    \end{subfigure}
    \hfill
    \begin{subfigure}[b]{0.48\textwidth}
        \centering
        \includegraphics[width=\textwidth]{plots/chc20_stromal_ecm_module_spatial_score.png}
        \caption{Stromal/ECM module score.}
        \label{fig:stromal_score}
    \end{subfigure}
    \vspace{0.4cm}

    \begin{subfigure}[b]{0.48\textwidth}
        \centering
        \includegraphics[width=\textwidth]{plots/chc20_combined_immune_stromal_niche_score.png}
        \caption{Combined microenvironment score.}
        \label{fig:combined_score}
    \end{subfigure}
    \hfill
    \begin{subfigure}[b]{0.48\textwidth}
        \centering
        \includegraphics[width=\textwidth]{plots/hcc4r_vs_chc20_combined_niche_score_distribution.png}
        \caption{Comparison of score distributions.}
        \label{fig:score_comparison}
    \end{subfigure}

    \caption{Spatial distribution and comparison of microenvironment module scores.}
    \label{fig:main_microenvironment_figure}
\end{figure}

Figures~\ref{fig:immune_score}, \ref{fig:stromal_score}, and \ref{fig:combined_score} show, respectively, the projections onto CHC20 spatial coordinates of the immune module score (derived from signature genes split by functional origin; $n = 15$ matched genes, encompassing core Treg/TAM molecules such as \textit{FOXP3}, \textit{IL2RA}, \textit{CTLA4}, \textit{TIGIT}, \textit{PDCD1}, \textit{LAG3}, \textit{CD163}, \textit{MRC1}, and \textit{CXCL12}), the stromal/ECM module score ($n = 23$ matched genes, encompassing CAF/ECM remodeling-related molecules such as \textit{FAP}, \textit{ACTA2}, \textit{COL1A1}, \textit{COL3A1}, \textit{SPP1}, \textit{MMP2}, \textit{LUM}, and \textit{DCN}), and the combined niche score obtained by Z-score standardizing and summing the two modules. The spatial patterns in all three panels are highly similar: high-scoring regions (yellow/orange) are concentrated within the same band-like area along the tissue margin in the upper-right portion of the section, whereas low-scoring regions (dark) are concentrated over a large area in the lower portion of the section, corresponding to the tumor-parenchyma-dense region in the CHC20 H\&E-stained section—indicating that the signature genes derived from screening the HCC4R section can effectively serve as a criterion for identifying distinct niches in other sections.

The violin plot in Fig.~\ref{fig:score_comparison} further compares the overall distribution of the combined score between the HCC4R and CHC20 sections. The Mann-Whitney U test comparing the combined niche score distributions between the two groups yields $p = 0.094$, which does not reach the conventional 0.05 significance threshold; the medians of the two violins (approximately $-0.5$ for HCC4R and $-0.2$ for CHC20), their interquartile ranges, and their overall shapes are highly overlapping, and although the upper and lower bounds (approximately $-2.9$ to $5.1$ for HCC4R and $-4.5$ to $7.4$ for CHC20) show a somewhat wider tail for CHC20, the bulk of the two distributions is essentially consistent.

Taken together, the high-scoring regions in CHC20 spatially overlap substantially with the immunosuppressive regions at the margins of tumor areas in its histological section, supporting the notion that the identified niche is a general histological feature of HCC rather than a phenomenon specific to a single section~\citep{llovet2021hepatocellular}. It is worth noting that CHC20 is a mixed hepatocellular-cholangiocarcinoma, and its immune microenvironment composition may differ biologically from that of classical HCC (HCC4R); against this background, the consistency of scores between the two sections is of particular significance.

\subsection{Independent prognostic validation in the TCGA-LIHC cohort}

\begin{figure}[h]
    \centering
    \begin{subfigure}[b]{0.48\textwidth}
        \centering
        \includegraphics[width=\textwidth]{plots/km_plot.png}
        \caption{Kaplan-Meier curves for the ssGSEA-high and ssGSEA-low groups.}
        \label{fig:km_plot}
    \end{subfigure}
    \hfill
    \begin{subfigure}[b]{0.48\textwidth}
        \centering
        \includegraphics[width=\textwidth]{plots/cox_forest_plot.png}
        \caption{Cox forest plot for the ssGSEA-high and ssGSEA-low groups.}
        \label{fig:cox_forest}
    \end{subfigure}
    \caption{Survival analysis results in the TCGA-LIHC cohort.}
    \label{fig:survival_result}
\end{figure}

The HCC4R spatial signature gene set was projected via ssGSEA onto patients in the TCGA-LIHC cohort, and patients were divided into High and Low groups based on the median ssGSEA score. Kaplan-Meier survival analysis (Fig.~\ref{fig:km_plot}) shows that the overall survival curves of the two groups are broadly similar, and the log-rank test does not reach statistical significance ($p = 0.79$), indicating that this spatial signature score fails to effectively discriminate overall survival differences among patients in the current TCGA-LIHC cohort.

Further multivariable Cox proportional hazards regression analysis (Fig.~\ref{fig:cox_forest}), after adjusting for clinical factors including age and tumor stage, shows a hazard ratio for the ssGSEA score close to 1 (HR = 0.92, 95\% CI: 0.36--2.36, $p = 0.862$), with a confidence interval spanning 1, indicating that this score does not exhibit independent prognostic value within the current model. By contrast, several tumor-stage variables show a markedly elevated hazard trend in the Cox model, suggesting that clinical stage remains an important determinant of overall survival in the TCGA-LIHC cohort.

This result indicates that, although the spatial niche signature reflects tissue composition, it is not an independent prognostic factor for bulk overall survival in TCGA-LIHC and shows little association with survival outcome.

\section{Discussion}

This study establishes a complete analytical framework spanning five stages—single-cell reference construction, spatial deconvolution, niche identification, mechanistic dissection, and cross-sample/clinical validation—achieving the goal of extracting spatial immune-regulatory principles from a single-cell transcriptomic reference atlas and systematically validating them in independent tissues and patient cohorts. Compared with conventional analytical paradigms that define tumor immunosuppressive status using prior marker gene sets, this approach offers three principal innovations. First, it employs an unsupervised Cellular Neighborhoods clustering strategy to discover immunosuppressive spatial structure directly from the data itself, rather than relying on pre-specified cell-type weights or thresholds. Second, to address the cell-dilution effect inherent to the Visium platform, it introduces the detection-fraction difference $\Delta$fraction as a compensatory dimension to conventional $\log_2$FC screening, effectively reducing the risk of false negatives arising from sparse expression and rendering the extraction of Treg-related marker gene features more sensitive. Third, it uses a rule-based spatial-domain labeling scheme as an independent semantic reference frame, cross-validating it against the data-driven Leiden clustering results at the methodological level, thereby substantially enhancing the robustness and interpretability of the immunosuppressive niche determination. This integrative strategy offers a reproducible methodological paradigm for combining single-cell-resolution reference atlases with spatial transcriptomic data, and its closed-loop ``discovery--mechanism--validation'' design has broader implications for microenvironmental spatial typing studies in other solid tumors.

Despite achieving highly consistent results at the histological level and across independent spatial sections, the TCGA-LIHC independent prognostic validation in this study did not meet expectations: neither the Kaplan-Meier survival analysis ($p = 0.79$) nor the multivariable Cox regression ($\text{HR} = 0.92$, $p = 0.862$) reached statistical significance, and neither the forest plot nor the KM curve yielded a risk-stratification signal of independent prognostic value from this spatial signature. This discrepancy may be attributable to several methodological limitations. First, the process of projecting ssGSEA scores onto bulk transcriptomes may severely dilute spatially specific signals; bulk RNA-seq measures average expression levels after homogenizing an entire tissue sample, discarding spatial localization information, such that locally high-signal regions that are truly driving the biology (e.g., the Treg-enriched invasive front or fibrotic stroma) may be masked by the much larger signal from normal hepatic parenchyma or bulk tumor tissue~\citep{svensson2018spatiallyvariable,gao2020loss}. Second, tumor tissue in the TCGA-LIHC cohort often exhibits varying degrees of necrosis, inflammation, and stromal composition that differ substantially among patients, and the bulk-level immune microenvironment composition is markedly influenced by sampling location~\citep{zhang2019landscape}; even within the same patient, image-guided biopsies from two different locations may yield substantially different intratumoral immune cell proportions, and this sampling bias may obscure any weak association between the bulk ssGSEA score and true overall survival. Third, along the causal chain from spatial transcriptomic features to overall survival, multiple intermediate steps exist; the immunosuppressive niche may be a local driver of short-term recurrence or treatment resistance following tumor resection without necessarily exerting a significant effect on overall survival~\citep{zhang2019landscape}. In addition, the clinical data in the TCGA-LIHC cohort are drawn from multiple contributing centers with heterogeneous surgical quality, adjuvant treatment regimens, and follow-up completeness, introducing further non-biological random noise~\citep{liu2018tcga}. Taken together, although the spatial niche signature precisely and reproducibly characterizes the immune spatial architecture, its clinical translation into a bulk overall-survival prognostic factor may require finer-grained endpoints (such as recurrence-free survival, RFS) or more balanced prospective sample-collection strategies~\citep{liu2018tcga}.

Beyond the non-significant survival analysis, this study also has several methodological limitations warranting future improvement. First, the immunosuppressive niche score depends on the choice of the neighborhood parameter $k$ and the radius multiplier; although sensitivity analysis has demonstrated overall robustness ($\rho > 0.90$, $\text{WJ} > 0.80$), the local spatial-boundary influence of extreme parameter variation cannot be entirely dismissed, and future work could consider a Bayesian hierarchical model that jointly infers posterior spatial patterns and neighborhood probability distributions, reducing dependence on a single fixed neighborhood definition~\citep{kleshchevnikov2022cell2location,gelman2013bayesian}. Second, although the ligand-receptor communication analysis introduces rank normalization to avoid the zero-inflation artifact of the raw product method, it does not fully distinguish genuine bidirectional communication from apparent high co-expression arising when ligand and receptor genes are expressed by different cell types within the same spot~\citep{cang2023commot}; this may require integration with a ligand-receptor gene co-localization probabilistic model. Third, the immunosuppressive gene module was screened based solely on upregulation, without accounting for immune evasion mechanisms mediated by the downregulation of suppressive molecules (such as loss of HLA genes)~\citep{montesion2021hla}. Fourth, the current analysis includes only a single HCC section, HCC4R, and a single mixed hepatocellular-cholangiocarcinoma section, CHC20, for cross-sample validation; the sample size remains limited, and future work should expand to multi-center, multi-cancer-type spatial transcriptomic cohorts.

Looking ahead, the analytical framework established in this study can be extended in three directions. First, registration and integration with multiplexed immunofluorescence (mIF) or MIBI/CODEX spatial proteomics data could further confirm, at the protein level, the spatial co-localization of key molecules within the immunosuppressive niche (the CXCL12--CXCR4 axis, the PD-1/PD-L1 immune checkpoint), providing molecular biomarkers for personalized efficacy prediction in immune checkpoint blockade (ICB) therapy~\citep{keren2018structured,zhu2021spark}. Second, as higher-resolution spatial sequencing technologies such as 10x Visium HD, Slide-seq, and Stereo-seq continue to mature, future work could re-evaluate, at near-single-cell or even subcellular resolution, whether the domain-definition scheme and feature gene signature established here remain applicable to fine-scale microenvironmental partitioning at the sub-spot level, further probing the fine topological structure of the tumor-immune interface~\citep{moses2022museum,chen2022spatia}. Third, the ``spatial niche signature'' could serve as an input feature for training deep learning models such as Transformers for label-free prediction of recurrence risk at surgical margins, or could be integrated with digital pathology to combine macroscopic H\&E image features with microscopic spatial transcriptomic signals in a multimodal framework, constructing a ``macro--micro'' cross-scale system for assessing tumor immune status~\citep{chen2016liver,shmatko2022artific}. The data-processing pipeline, feature-selection criteria, and sensitivity-analysis framework accumulated in this study may provide a foundational methodological reference for these extensions.

% ============================================================
%  Part V (Appendix): Implementation Details
% ============================================================

\newpage
\appendix
\renewcommand{\thesection}{\Alph{section}}

\section{List of ligand-receptor communication pairs}
\label{app:lr_table}

\begin{table}[H]
\centering
\caption{Ligand-receptor communication pairs included in the analysis.}
\label{tab:lr_pairs}
\begin{tabularx}{\textwidth}{lllX}
\toprule
\textbf{Ligand} & \textbf{Receptor} & \textbf{Pair} & \textbf{Biological significance} \\
\midrule
CXCL12 & CXCR4   & CXCL12--CXCR4  & CAF-secreted; recruits Treg/MDSC~\citep{domanska2013targeting} \\
TGFB1  & TGFBR1  & TGFB1--TGFBR1  & Immunosuppressive cytokine; co-secreted by CAF and Treg~\citep{mariathasan2018tgfb} \\
PDCD1  & CD274   & PD-1--PD-L1    & Classical immune checkpoint axis~\citep{sharma2023immunosuppressive} \\
TIGIT  & NECTIN2 & TIGIT--NECTIN2 & Treg inhibitory checkpoint~\citep{sakaguchi2020regulatory} \\
LAG3   & HLA-DRA & LAG3--MHC-II   & Associated with T-cell exhaustion~\citep{sharma2023immunosuppressive} \\
CCL22  & CCR4    & CCL22--CCR4    & TAM-mediated Treg recruitment~\citep{curiel2004specific} \\
IL10   & IL10RA  & IL10--IL10RA   & Treg immunosuppressive cytokine axis~\citep{noy2014tumor} \\
SPP1   & CD44    & SPP1--CD44     & TAM-secreted osteopontin~\citep{biswas2010macrophage} \\
MIF    & CD74    & MIF--CD74      & Macrophage migration inhibitory factor~\citep{noy2014tumor} \\
VEGFA  & KDR     & VEGFA--KDR     & Angiogenesis; CAF-related~\citep{kalluri2016biology} \\
LGALS9 & HAVCR2  & Galectin9--TIM-3 & Inhibitory checkpoint~\citep{sharma2023immunosuppressive} \\
CCL2   & CCR2    & CCL2--CCR2     & Macrophage recruitment pathway~\citep{biswas2010macrophage} \\
\bottomrule
\end{tabularx}
\end{table}

\section{Definitions of neighborhood-parameter sensitivity metrics}
\label{app:sensitivity}

The neighborhood-parameter sensitivity analysis described in Section 2.4 of the main text employs two complementary continuous metrics to jointly evaluate the consistency of the composite niche score before and after perturbing the $k$-NN neighbor count and the radius multiplier.

The first metric is the Spearman rank correlation coefficient $\rho$, defined as
\begin{equation}
\rho = 1 - \frac{6 \sum_i d_i^2}{n(n^2-1)}
\label{eq:spearman}
\end{equation}
where $d_i$ is the difference in rank between the reference score and the perturbed score at spot $i$, and $n$ is the total number of spots. $\rho$ directly measures the consistency of the relative ordering of all spots, is unaffected by any truncation threshold, and ranges over $[-1, 1]$; values closer to 1 indicate that the parameter perturbation has almost no effect on spot ranking.

The second metric is the continuous weighted Jaccard similarity (Weighted Jaccard), defined as
\begin{equation}
\text{WJ}(\mathbf{a}, \mathbf{b}) = \frac{\sum_i \min(a_i, b_i)}{\sum_i \max(a_i, b_i)}
\label{eq:weighted_jaccard}
\end{equation}
where $\mathbf{a}$ and $\mathbf{b}$ are, respectively, the reference score vector and the perturbed score vector after min-max normalization to $[0,1]$. Compared with the traditional Jaccard index based on binary sets, this continuous form naturally assigns greater weight to high-scoring regions (i.e., candidate niche regions), compensating for the fact that Spearman rank correlation treats the top and bottom of the distribution equally; it ranges over $[0, 1]$, with values closer to 1 indicating a high degree of spatial overlap among high-scoring regions.

\section{Auxiliary scoring formulae}
\label{app:scores}

The auxiliary scoring system is defined as follows, where $z(\cdot)$ denotes Z-score standardization across spots:

\begin{align}
S_{\text{Treg-like},i} &= z(\hat{p}_{\text{Treg},i}) + z(\hat{s}_{\text{gene},i}) \\
S_{\text{immune-stroma},i} &= z(\hat{p}_{\text{Treg},i}) + z(\hat{p}_{\text{T/NK},i}) + z(\hat{p}_{\text{Myeloid},i}) + z(\hat{p}_{\text{Fibroblast},i}) \\
S_{\text{niche},i} &= z(\hat{p}_{\text{Treg},i}) + z(\hat{p}_{\text{Myeloid},i}) + z(\hat{p}_{\text{Fibroblast},i}) + z(\hat{s}_{\text{gene},i}) + z(\bar{n}_{\text{Hep},i})
\end{align}

where $\hat{p}$ is the Cell2location posterior abundance normalized to a proportion, $\hat{s}_{\text{gene}}$ is the immunosuppressive gene module score, and $\bar{n}_{\text{Hep}}$ is the mean hepatocyte proportion within a spot's neighborhood (used to capture the tumor-stroma interface effect). Dual-module scores (used for CHC20 validation):

\begin{align}
\text{immune\_score}_i &= \frac{1}{|G_{\text{immune}}|}\sum_{g \in G_{\text{immune}}} E_{ig} \\
\text{stromal\_score}_i &= \frac{1}{|G_{\text{stromal}}|}\sum_{g \in G_{\text{stromal}}} E_{ig} \\
\text{combined\_niche\_score}_i &= z(\text{immune\_score}_i) + z(\text{stromal\_score}_i)
\end{align}

$E_{ig}$ is the log1p-normalized (CP10K) expression level; the immune module comprises 21 genes including \textit{FOXP3}, \textit{IL2RA}, \textit{CTLA4}, and \textit{CD163}, and the stromal/ECM module comprises 21 genes including \textit{FAP}, \textit{ACTA2}, \textit{COL1A1}, and \textit{SPP1}.

\section{Core code implementation}
\label{app:code}

\subsection{Cell2location deconvolution}

\begin{lstlisting}[caption={Core code for RegressionModel training and spatial modeling}]
# Stage 1: train the cell-type reference signature matrix
from cell2location.models import RegressionModel
RegressionModel.setup_anndata(adata_sc, labels_key='final_celltype')
mod = RegressionModel(adata_sc)
mod.train(max_epochs=250, batch_size=1024, use_gpu=True,
          early_stopping=True, early_stopping_patience=30)
cell_state_df = mod.export_posterior(adata_sc,
    use_cols=mod.feature_in_columns)['means_per_cluster_mu_fg']

# Stage 2: spatial modeling
import cell2location
cell2location.models.Cell2location.setup_anndata(adata_vis)
mod = cell2location.models.Cell2location(adata_vis,
    cell_state_df=cell_state_df,
    N_cells_per_location=30, detection_alpha=20)
mod.train(max_epochs=1000, batch_size=1000, use_gpu=True)
adata_vis = mod.export_posterior(adata_vis,
    sample_kwargs={'num_samples': 1000, 'batch_size': 1000})
# Results stored in adata_vis.obsm['means_cell_abundance_w_sf']
\end{lstlisting}

\subsection{Neighborhood composition clustering and niche annotation}

\begin{lstlisting}[caption={Neighborhood composition vector computation, Leiden clustering, and composite-rank annotation}]
from scipy.spatial import KDTree
import scanpy as sc

# Construct the kNN adjacency graph and compute the neighborhood composition vector
tree = KDTree(coords)
_, indices = tree.query(coords, k=k + 1)  # exclude self
neighbors = [indices[i, 1:] for i in range(len(coords))]
arr = proportions.to_numpy()
neighborhood_comp = np.array(
    [arr[idx].mean(axis=0) if len(idx) else arr[i]
     for i, idx in enumerate(neighbors)])

# Leiden clustering (directly in composition-vector feature space, no PCA needed)
adata_nc = ad.AnnData(neighborhood_comp.astype(float))
sc.pp.neighbors(adata_nc, n_neighbors=n_neighbors, use_rep='X')
sc.tl.leiden(adata_nc, resolution=0.5, key_added='nc_cluster')

# Composite-rank annotation of the immunosuppressive niche
stats = df.groupby('nc_cluster')[
    [treg_col, myeloid_col, fibroblast_col, gene_score_col]].mean()
for col in stats.columns:
    stats[f'rank_{col}'] = stats[col].rank(ascending=False)
stats['combined_rank'] = stats.filter(like='rank_').sum(axis=1)
niche_id = stats['combined_rank'].idxmin()
\end{lstlisting}

\subsection{Two-layer feature gene screening}

\begin{lstlisting}[caption={Layer 1 differential testing and Layer 3 Gini index computation}]
# Layer 1: block-level Mann-Whitney U + AUC
from scipy.stats import mannwhitneyu
# Spatial grid aggregation (reduces spatial autocorrelation)
gx = np.floor(df['spatial_x'].values / block_size).astype(int)
gy = np.floor(df['spatial_y'].values / block_size).astype(int)
block_ids = [f'{x}_{y}' for x, y in zip(gx, gy)]
# Compute AUC = U / (n_high * n_low)
u_stat, pval = mannwhitneyu(high_vals, low_vals, alternative='greater')
auc = u_stat / (len(high_vals) * len(low_vals))

# Layer 3: Gini index for filtering focally sparse genes
def gini_index(x):
    x = np.sort(x.ravel())
    n = len(x)
    cumx = np.cumsum(x)
    return (n + 1 - 2 * cumx.sum() / cumx[-1]) / n if cumx[-1] > 0 else 0
\end{lstlisting}

\subsection{Dual-module scoring and ssGSEA}

\begin{lstlisting}[caption={Dual-module score computation (Python) and ssGSEA scoring (R)}]
# Python: dual-module scoring
def compute_dual_module_scores(adata, immune_genes, stromal_genes):
    expr = _expression_frame(adata)  # CP10K + log1p
    immune_score = expr[[g for g in immune_genes
                          if g in expr.columns]].mean(axis=1)
    stromal_score = expr[[g for g in stromal_genes
                           if g in expr.columns]].mean(axis=1)
    def _z(s): return (s - s.mean()) / s.std() if s.std() > 0 else s * 0
    return immune_score, stromal_score, _z(immune_score) + _z(stromal_score)
\end{lstlisting}

\begin{lstlisting}[language=R, caption={ssGSEA and multivariable Cox regression (R)}]
library(GSVA); library(survival); library(survminer)
gsva_scores <- gsva(expr_matrix, list(sig = signature_genes),
    method='ssgsea', kcdf='Gaussian', abs.ranking=TRUE, min.sz=5)
cox_model <- coxph(
    Surv(OS_time, OS_event) ~ ssgsea_score + age + stage,
    data = clinical_data)
summary(cox_model)
\end{lstlisting}

% ============================================================
% References
% ============================================================
\bibliographystyle{unsrt}
\bibliography{report_0701}

\end{document}
